\problem{3.1}{Square-integrable functions form a vector space and inner-product space}
Let $L^2(a,b)$ denote equivalence classes of measurable complex functions
with finite $\|f\|^2=\int_a^b|f|^2\dd x$.
For $f,g\in L^2$ and complex $\alpha,\beta$, use
$|u+v|^2\le2|u|^2+2|v|^2$ to obtain
\begin{equation}
\|\alpha f+\beta g\|^2\le
2|\alpha|^2\|f\|^2+2|\beta|^2\|g\|^2<\infty.
\end{equation}
Thus the set is closed under vector addition and scalar multiplication;
the other vector-space axioms hold pointwise almost everywhere.
The \emph{normalized} subset is not a vector space: it does
not contain zero and, for normalized $f$, the multiple $2f$
has norm two.
For $\braket f g=\int f^*g$, linearity in the second slot gives
$\braket f{\alpha g+\beta h}=\alpha\braket fg+
\beta\braket fh$. Conjugate symmetry follows by complex
conjugation under the integral; positivity is
$\braket ff=\int|f|^2\ge0$, with equality exactly
for the zero equivalence class. The Cauchy--Schwarz
inequality assures finite overlaps.
\answer{$L^2(a,b)$ is an inner-product vector space and is complete
as a normed space; the unit sphere of normalized states is not
a linear subspace.}

\problem{3.2}{Integrability and the unbounded derivative operator}
For $f(x)=x^\nu$, $0<x<1$, square integrability requires
\begin{equation}
\int_0^1|f|^2\dd x
=\int_0^1x^{2\nu}\dd x
=\frac{1}{2\nu+1}<\infty
\quad\Longleftrightarrow\quad\boxed{\nu>-\frac12}.
\end{equation}
For $\nu=1/2$, $f=\sqrt x$ and $xf=x^{3/2}$
are in $L^2(0,1)$, with squared norms $1/2$ and $1/4$.
But $f'=1/(2\sqrt x)$ has
$\int_0^1|f'|^2\dd x=\frac14\int_0^1\dd x/x=\infty$.
\answer{Differentiation does not map every $L^2$ function
back into $L^2$; the derivative operator requires a domain.}

\problem{3.3}{Equivalence of Hermiticity tests via polarization}
Suppose $\braket h{Qh}=\braket{Qh}h$ for all
$h$ in a complex linear domain.
Set $B(f,g)=\braket f{Qg}-\braket{Qf}g$.
Then $B(h,h)=0$, while $B$ is conjugate-linear
in its first entry and linear in its second.
Expanding $B(f+g,f+g)=0$ gives
\begin{equation}
B(f,g)+B(g,f)=0.
\end{equation}
Expanding $B(f+\ii g,f+\ii g)=0$ instead gives
$\ii B(f,g)-\ii B(g,f)=0$. Together they imply
$B(f,g)=0$ for arbitrary $f,g$:
\begin{equation}
\boxed{\braket f{Qg}=\braket{Qf}g}.
\end{equation}
The converse follows by setting $f=g$.
\qnote{This proves the algebraic symmetry identity;
self-adjointness additionally requires equality of adjoint domains.}

\problem{3.4}{Sums, scalar multiples and products of Hermitian operators}
On an appropriate common dense domain,
$(A+B)^\dagger=A^\dagger+B^\dagger$;
thus the sum of symmetric operators is symmetric.
For a scalar $\alpha$ and nonzero Hermitian $Q$,
$(\alpha Q)^\dagger=\alpha^*Q$, so
$\alpha Q$ is Hermitian precisely when $\alpha$ is real.
For Hermitian $A,B$, $(AB)^\dagger=BA$;
$AB$ is Hermitian when $AB=BA$ (with the
necessary shared operator-domain conditions).
Position is multiplication by real $x$:
$\int f^*xg=\int(xf)^*g$.
For $H=-\hbar^2\partial_x^2/(2m)+V(x)$,
two integrations by parts give
\begin{equation}
\braket f{Hg}-\braket{Hf}g
=-\frac{\hbar^2}{2m}
\left[f^*g'-f'^*g\right]_{\partial I},
\end{equation}
where $V$ is real. The surface term vanishes
for suitable boundary conditions, proving formal symmetry.
\qnote{Sums and products of unbounded self-adjoint
operators require domain care; the finite-dimensional
matrix rules are not automatically sufficient.}

\problem{3.5}{Adjoint rules and oscillator ladder operators}
Integration by parts shows, when the boundary
term vanishes, $(\partial_x)^\dagger=-\partial_x$.
Multiplication by $x$ satisfies $x^\dagger=x$,
whereas $(\ii I)^\dagger=-\ii I$. For operators
on a common compatible domain,
\begin{equation}
\boxed{(QR)^\dagger=R^\dagger Q^\dagger,\quad
(Q+R)^\dagger=Q^\dagger+R^\dagger,\quad
(cQ)^\dagger=c^*Q^\dagger}.
\end{equation}
These follow directly from applying the adjoint
definition twice: $\braket f{QRg}
=\braket{Q^\dagger f}{Rg}
=\braket{R^\dagger Q^\dagger f}g$.
The oscillator raising operator
$a_+=(m\omega x-\ii p)/\sqrt{2m\hbar\omega}$
therefore has adjoint
\begin{equation}
\boxed{a_+^\dagger=
\frac{m\omega x+\ii p}{\sqrt{2m\hbar\omega}}=a_-}.
\end{equation}

\problem{3.6}{Angular second derivative, eigenvalues and degeneracy}
On the periodic domain $f(\phi+2\pi)=f(\phi)$,
$f'(\phi+2\pi)=f'(\phi)$, integrate twice:
\begin{equation}
\int_0^{2\pi}f^*g''\dd\phi=
[f^*g'-f'^*g]_0^{2\pi}
+\int_0^{2\pi}f''{}^*g\dd\phi
=\int_0^{2\pi}f''{}^*g\dd\phi.
\end{equation}
Hence $Q=\partial_\phi^2$ is symmetric
(and self-adjoint on the usual periodic $H^2$ domain).
The normalized eigenfunctions are
\begin{equation}
\boxed{u_n(\phi)=\frac{\ee^{\ii n\phi}}{\sqrt{2\pi}}},
\quad\boxed{Q u_n=-n^2u_n},\qquad n\in\mathbb Z.
\end{equation}
The spectrum is $0,-1,-4,-9,\ldots$; $0$ is simple,
each $-n^2$ ($n\ge1$) is doubly degenerate,
with independent $\ee^{\pm\ii n\phi}$.

\problem{3.7}{Degenerate eigenspaces and orthogonal combinations}
If $Qf=qf$ and $Qg=qg$, linearity gives
$Q(\alpha f+\beta g)=q(\alpha f+\beta g)$.
For $Q=\partial_x^2$ on $(-1,1)$,
$f=\ee^x$ and $g=\ee^{-x}$ both have eigenvalue $1$.
Choose their even/odd combinations
\begin{equation}
u(x)=\cosh x=\frac{\ee^x+\ee^{-x}}2,\qquad
v(x)=\sinh x=\frac{\ee^x-\ee^{-x}}2.
\end{equation}
Both satisfy $u''=u$, $v''=v$, and
$\int_{-1}^1u(x)v(x)\dd x=0$
because the integrand is odd.
\qnote{These are orthogonal eigenfunctions of the
differential expression. A self-adjoint boundary
domain must be specified before interpreting the
degeneracy as a physical observable's spectrum.}

\problem{3.8}{Orthogonality of angular eigenfunctions}
For the angular momentum operator of Example 3.1,
$L_z=-\ii\hbar\partial_\phi$,
the periodic eigenfunctions
$u_n(\phi)=\ee^{\ii n\phi}/\sqrt{2\pi}$
obey $L_zu_n=n\hbar u_n$: the eigenvalues are real.
Their overlaps are
\begin{equation}
\braket{u_m}{u_n}=
\frac1{2\pi}\int_0^{2\pi}
\ee^{\ii(n-m)\phi}\dd\phi
=\delta_{mn}.
\end{equation}
For Problem 3.6, the same basis diagonalizes
$\partial_\phi^2$ with eigenvalues $-n^2$.
Distinct eigenvalues correspond to distinct
$|n|$, hence orthogonality follows from the
same integral. For $n$ and $-n$, orthogonality
still holds despite the degeneracy (for $n\ne0$).

\problem{3.9}{Discrete, continuous and mixed energy spectra}
A particle in a one-dimensional infinite well
has only the discrete energies
$E_n=n^2\pi^2\hbar^2/(2ma^2)$.
A potential step has only scattering states in
the customary no-bound-state configuration,
and hence a continuous energy spectrum.
An attractive delta well has one bound level
$E_b=-m\alpha^2/(2\hbar^2)$ and positive-energy
scattering states; its spectrum is mixed.
\qnote{Other examples may be used if the
potential and boundary conditions are stated:
the type of spectrum depends on the self-adjoint
Hamiltonian, not merely its differential formula.}

\problem{3.10}{Is the ground well state a momentum eigenstate?}
For $0<x<a$, the normalized ground state is
$\psi_1(x)=\sqrt{2/a}\sin(\pi x/a)$.
Its derivative gives
\begin{equation}
\hat p\psi_1=-\ii\hbar\frac{\pi}{a}
\sqrt{\frac2a}\cos(\pi x/a),
\end{equation}
which is not a constant multiple of $\psi_1$.
Therefore it is not a momentum eigenstate.
An energy eigenstate in a box is a standing wave,
not a single traveling momentum wave.
Although $\avg p=0$ and
$\avg{p^2}=\pi^2\hbar^2/a^2$, neither result
makes the state an eigenvector of $\hat p$.
\qnote{There is a subtle distinction between the
formal derivative and the domain of a self-adjoint
momentum operator on a finite interval,
explored explicitly in Problem 3.48.}
